%------------------------------------------------------------------------------------
% ЗАКЛЮЧЕНИЕ
%------------------------------------------------------------------------------------
\newpage
\begin{center}
  \textbf{\large ЗАКЛЮЧЕНИЕ}
\end{center}
\addcontentsline{toc}{chapter}{ЗАКЛЮЧЕНИЕ}

В ходе выполнения дипломной работы была достигнута поставленная цель: 
исследовано влияние доплеровского масштабирования спектра на 
производительность системы спутниковой связи на примере стандарта 5G NR NTN. 
Решение всех сформулированных задач позволило получить следующие основные результаты.

\begin{enumerate}
    \item \textbf{Проведён обзор литературы.} 
    Рассмотрены актуальные работы, посвящённые развитию сетей NTN, 
    стандартам 5G NR NTN в релизах 3GPP, 
    а также моделям каналов связи. 

    \item \textbf{Получена теоретическая оценка межканальной интерференции (ICI).} 
    Показано, что в широкополосных системах (OFDM) 
    индивидуальные сдвиги поднесущих приводят к ICI, причём величина ICI пропорциональна количеству поднесущих
    и относитель\-ной скорости движения приемника/передатчика.
    Получено выражение для уровня SINR с учетом интерференции между поднесущими.

    \item \textbf{Проведено численное моделирование.} 
    Разработана имитационная модель OFDM-системы, позволяющая оценивать вероятность битовой ошибки (BER) 
    при различных параметрах: количестве поднесущих, видах модуляции и относительных скоростях.

    \item \textbf{Оценено влияние эффекта масштабирования спектра на производитель\-ность.}
    Установлено, что доплеровское масштабирование приводит к существенному ухудшению помехоустойчивости в сценариях, 
    характерных для низкоорбитальных спутников. Для широкополосных сигналов ($N=2048$) и высокоуровневых модуляций (КАМ-64, КАМ-256) 
    потери в SNR достигают 7 дБ, а при КАМ-256 наблюдается плато ошибок, делающее демодуляцию невозможной. 
    Применение частотного планирования 5G NR NTN позволяет сократить потери до 4 дБ для КАМ-64 и сделать влияние незначительным 
    для КАМ-16.
\end{enumerate}

\end{enumerate}